\documentclass[11pt]{article}
% ============================================================================
% preamble.tex — Estilo común de los manuales de práctica SISELA (FIME / UANL)
%
% Reproduce (de forma aproximada) el estilo de los PDF originales
% (~/Descargas/manuales/Manuales/pN.pdf), de los que no existe fuente .tex.
% Compilar con:  pdflatex  (o  latexmk -pdf pN.tex)
% ============================================================================

\usepackage[utf8]{inputenc}
\usepackage[T1]{fontenc}
\usepackage[spanish,es-tabla]{babel}
\usepackage[a4paper,margin=2.5cm,headheight=15pt]{geometry}

% --- Tipografía: serif para el cuerpo, sans para los títulos ---
\usepackage{libertinus}
\usepackage[scaled=0.92]{helvet}   % sans para encabezados
\usepackage{microtype}

% --- Matemáticas ---
\usepackage{amsmath,amssymb}
\usepackage{siunitx}
\sisetup{detect-all, per-mode=symbol}

% --- Tablas y figuras ---
\usepackage{booktabs}
\usepackage{array}
\usepackage{multirow}
\usepackage{caption}
\captionsetup{font=small,labelfont=bf,labelsep=period}
\usepackage{graphicx}
\usepackage{float}

% --- Diagramas ---
\usepackage{tikz}
\usetikzlibrary{shapes.geometric,arrows.meta,positioning,calc}
\tikzset{
  block/.style   = {draw, rectangle, minimum width=2.6cm, minimum height=0.9cm,
                    align=center, font=\footnotesize},
  term/.style    = {draw, rounded corners, minimum width=2.2cm, minimum height=0.8cm,
                    align=center, font=\footnotesize, fill=black!5},
  decision/.style= {draw, diamond, aspect=2, align=center, font=\scriptsize,
                    inner sep=1pt},
  flow/.style    = {-{Latex[length=2mm]}, thick},
}

% --- Código ---
\usepackage{xcolor}
\definecolor{codebg}{gray}{0.96}
\definecolor{codekw}{rgb}{0.13,0.29,0.53}
\definecolor{codecm}{rgb}{0.40,0.45,0.40}
\usepackage{listings}
\lstdefinestyle{sisela}{
  backgroundcolor=\color{codebg}, basicstyle=\ttfamily\footnotesize,
  keywordstyle=\color{codekw}\bfseries, commentstyle=\color{codecm}\itshape,
  stringstyle=\color{codekw}, numbers=left, numberstyle=\tiny\color{gray},
  numbersep=6pt, frame=none, breaklines=true, showstringspaces=false,
  columns=fullflexible, xleftmargin=1.8em, aboveskip=1em, belowskip=1em,
  literate={á}{{\'a}}1 {é}{{\'e}}1 {í}{{\'i}}1 {ó}{{\'o}}1 {ú}{{\'u}}1
           {ñ}{{\~n}}1 {Á}{{\'A}}1 {°}{{\textdegree}}1 {µ}{{$\mu$}}1
           {✓}{{\checkmark}}1,
}
\lstset{style=sisela}
\newcommand{\code}[1]{\texttt{\small #1}}

% --- Listas ---
\usepackage{enumitem}
\setlist{nosep, leftmargin=1.6em}

% --- Encabezado / pie de página ---
\usepackage{fancyhdr}
\usepackage{lastpage}
\pagestyle{fancy}
\fancyhf{}
\fancyhead[L]{\footnotesize\sffamily Sistemas Eléctricos de Aeronaves — Práctica \practnum}
\fancyhead[R]{\footnotesize\sffamily FIME / UANL}
\fancyfoot[C]{\footnotesize P\'agina \thepage\ de \pageref{LastPage}}
\renewcommand{\headrulewidth}{0.4pt}
\renewcommand{\footrulewidth}{0.4pt}

% --- Títulos de sección en sans-serif ---
\usepackage{titlesec}
\titleformat{\section}{\normalfont\Large\sffamily\bfseries}{\thesection}{0.7em}{}
\titleformat{\subsection}{\normalfont\large\sffamily\bfseries}{\thesubsection}{0.6em}{}
\titleformat{\subsubsection}{\normalfont\normalsize\sffamily\bfseries}{\thesubsubsection}{0.5em}{}

% --- hyperref al final ---
\usepackage[hidelinks,pdfusetitle]{hyperref}

% ============================================================================
% Comandos de portada y cajas
% ============================================================================
\newcommand{\practnum}{N}          % lo redefine cada pN.tex
\newcommand{\practtitle}{}         % idem

\newcommand{\portada}{%
  \thispagestyle{empty}
  \begin{center}
    {\large Universidad Autónoma de Nuevo León}\\[2pt]
    {\normalsize Facultad de Ingeniería Mecánica y Eléctrica}\\[2.2cm]
    \rule{\linewidth}{0.4pt}\\[0.4cm]
    {\Huge\sffamily\bfseries Práctica \practnum}\\[0.5cm]
    {\large \practtitle}\\[0.3cm]
    \rule{\linewidth}{0.4pt}\\[2.5cm]
    {\normalsize Laboratorio de Sistemas Electrónicos de Aeronaves\\(SISELA)}\\[1.5cm]
    {\normalsize Ciclo Escolar 2026}\\[3cm]
  \end{center}
}

\newenvironment{resumen}{\par\noindent\rule{\linewidth}{0.4pt}\par\smallskip
  \small\itshape\noindent\textbf{Resumen Ejecutivo} ---\ }{\par\smallskip
  \noindent\rule{\linewidth}{0.4pt}\par}

\newenvironment{nota}[1][Nota]{\par\smallskip\noindent\rule{\linewidth}{0.4pt}\par\smallskip
  \small\itshape\noindent\textbf{#1} ---\ }{\par\smallskip
  \noindent\rule{\linewidth}{0.4pt}\par\smallskip}

\renewcommand{\practnum}{7}
\renewcommand{\practtitle}{Control de Motores a Pasos\\
Posicionamiento Angular Discreto — Énfasis Aeronáutico}

\begin{document}
\portada

\begin{resumen}
Los motores a pasos (\emph{stepper motors}) son dispositivos electromecánicos que
convierten pulsos eléctricos en desplazamientos angulares discretos. A diferencia de los
motores DC convencionales o los servomotores, el motor a pasos avanza una fracción de
vuelta precisa por cada pulso recibido, permitiendo un control de posición en lazo
abierto con alta repetibilidad. En el ámbito aeronáutico, los motores a pasos encuentran
su nicho principal en la instrumentación analógica de cabina (\emph{glass-analog
hybrid}): las agujas de indicadores como altímetros, velocímetros, indicadores de RPM,
temperatura de gases de escape (EGT) y presión de combustible son accionadas por pequeños
motores que reciben comandos de posicionamiento desde el bus de datos ARINC 429. Esta
práctica cubre el motor unipolar 28BYJ-48 con su driver ULN2003, el motor bipolar
NEMA 17 con el driver A4988/DRV8825, homing con fin de carrera, y posicionamiento
absoluto vía ARINC 429 Label 204.
\end{resumen}

\tableofcontents
\newpage

% ============================================================================
\section{Introducción}

Los motores a pasos (\emph{stepper motors}) son dispositivos electromecánicos que
convierten pulsos eléctricos en desplazamientos angulares discretos. A diferencia de los
motores DC convencionales que giran continuamente al aplicar voltaje, o de los
servomotores que requieren un lazo de retroalimentación, el motor a pasos avanza una
fracción de vuelta precisa por cada pulso recibido, permitiendo un control de posición en
lazo abierto con alta repetibilidad.

En el ámbito aeronáutico, los motores a pasos encuentran su nicho principal en la
\textbf{instrumentación analógica de cabina} (\emph{glass-analog hybrid}). Las agujas de
indicadores como altímetros, velocímetros, indicadores de RPM, temperatura de gases de
escape (EGT) y presión de combustible son accionadas por pequeños motores que reciben
comandos de posicionamiento desde el bus de datos ARINC 429. El procesador del
instrumento decodifica la palabra ARINC 429, calcula la posición angular correspondiente
y envía los pasos necesarios al \emph{driver} para actualizar la carátula con alta
resolución y repetibilidad.

Adicionalmente, los motores a pasos se emplean en posicionamiento de antenas pequeñas,
control de válvulas de bajo torque, equipos de prueba y calibración, y en sistemas de
prueba de vuelo (\emph{flight test instrumentation}).

En las prácticas anteriores se cubrió el control con servomotores (P5, posición por PWM)
y la conmutación de potencia con transistores BJT (P6, diseño de interfaces
\emph{low-side} y \emph{high-side}). Esta práctica da el siguiente paso: controlar la
\textbf{posición angular discreta} de una carga mecánica usando secuencias de conmutación
de bobinas, sin necesidad de retroalimentación analógica.

El estudiante trabajará con el motor unipolar \textbf{28BYJ-48} y su \emph{driver}
\textbf{ULN2003} como caso de estudio principal. Esta combinación es ideal para
prototipos educativos por su bajo costo, alimentación a 5\,V y control directo desde
pines GPIO. Como alternativa, se presenta el motor bipolar \textbf{NEMA 17} con el
\emph{driver} A4988/DRV8825, comparando ambas arquitecturas.

% ============================================================================
\section{Objetivos de Aprendizaje}
\begin{enumerate}
  \item Identificar las diferencias constructivas y operativas entre motores a pasos
        unipolares (28BYJ-48, 5 cables) y bipolares (NEMA 17, 4 cables), y sus
        respectivos \emph{drivers} de potencia.
  \item Diseñar e implementar las secuencias lógicas de conmutación \emph{full-step}
        (4 fases) y \emph{half-step} (8 fases) en software, controlando los pines
        IN1--IN4 del ULN2003.
  \item Calcular la resolución angular del sistema considerando la reducción mecánica
        interna del motor: $\text{SPR}=360\si{\degree}/\theta_\text{paso}\times
        n_\text{reductor}$ y relacionar RPM con el intervalo de tiempo entre pasos.
  \item Implementar un sistema de referenciación (\emph{homing}) utilizando un sensor de
        fin de carrera para establecer la posición cero del instrumento.
  \item Integrar la recepción de tramas ARINC 429 (Label 204) para el posicionamiento
        absoluto de un indicador, incluyendo validación de paridad y cálculo de
        trayectoria mínima.
  \item Comprender el funcionamiento del ULN2003 como arreglo de pares Darlington con
        diodos de protección internos, y su relación con la conmutación de bobinas
        inductivas.
  \item Comparar el control basado en secuencias (ULN2003) con el control STEP/DIR
        (A4988), evaluando torque, resolución, complejidad y aplicabilidad aeronáutica.
  \item Implementar un control interactivo con menú serial: \emph{jog} manual,
        movimiento por pasos, barrido y \emph{homing}, verificando la operación con
        osciloscopio y midiéndola con instrumentos.
  \item Documentar el trabajo en formato AIAA: teoría, cálculos, diagramas, código
        fuente, resultados experimentales y análisis estadístico de precisión.
  \item Contextualizar los motores a pasos en aplicaciones aeronáuticas: instrumentación
        de cabina, posicionamiento de antenas, control de válvulas y sistemas de
        prueba, considerando normativas DO-178C y DO-160G.
\end{enumerate}

% ============================================================================
\section{Fundamentos Teóricos}

\subsection{Principio de funcionamiento del motor a pasos}
Un motor a pasos es un dispositivo electromecánico en el que el rotor (generalmente de
imán permanente o reluctancia variable) se alinea con el campo magnético generado por las
bobinas del estátor. Al energizar las bobinas en una secuencia específica, el campo
magnético ``rota'' discretamente y el rotor lo sigue en incrementos angulares fijos
denominados \textbf{pasos}. Los parámetros fundamentales son:
\begin{itemize}
  \item \textbf{Ángulo de paso} ($\theta_\text{paso}$): incremento angular por cada paso
        del motor. Valores comunes: $1.8\si{\degree}$ (NEMA 17), $5.625\si{\degree}$
        (28BYJ-48 antes de reducción).
  \item \textbf{Pasos por revolución}: $\text{SPR}=360\si{\degree}/\theta_\text{paso}$.
        Para el eje interno del 28BYJ-48: 64 pasos.
  \item \textbf{Torque de mantenimiento} (\emph{holding torque}): torque máximo cuando
        las bobinas están energizadas y el rotor está detenido.
  \item \textbf{Torque de desprendimiento} (\emph{pull-out torque}): torque máximo a una
        velocidad dada sin perder pasos.
\end{itemize}

\begin{nota}[Lazo abierto vs lazo cerrado]
El control en lazo abierto es viable siempre que no se exceda el torque de
desprendimiento. Si la carga mecánica supera la capacidad del motor, se pierden pasos sin
que el controlador lo detecte. En instrumentación aeronáutica, un procedimiento de
\emph{homing} periódico compensa esta limitación.
\end{nota}

\subsection{Tipos de motores a pasos}
\subsubsection{Motor unipolar (28BYJ-48)}
Cada bobina tiene una toma central (\emph{center tap}) conectada a $V_{CC}$. La corriente
fluye siempre en la misma dirección en cada sección de bobina, simplificando el
\emph{driver} a transistores individuales (sin necesidad de puente H).
\begin{itemize}
  \item \textbf{Cables}: 5 (rojo = común $V_{CC}$; azul, rosa, amarillo, naranja = fases
        A--D).
  \item \textbf{Voltaje}: 5\,V DC.
  \item \textbf{Corriente por fase}: $\sim$\SIrange{150}{200}{\milli\ampere}.
  \item \textbf{Reducción interna}: $\approx 64{:}1$ (real $\approx 63.68{:}1$).
  \item \textbf{Pasos por revolución del eje de salida}: 2048 (\emph{full-step}) / 4096
        (\emph{half-step}).
  \item \textbf{Torque}: $\sim$\SIrange{0.3}{0.5}{\kilo\gram\centi\meter} (bajo, adecuado
        para agujas e indicadores).
  \item \textbf{Velocidad máxima}: $\sim$\SIrange{10}{15}{RPM} (limitado por la caja
        reductora).
\end{itemize}

\subsubsection{Motor bipolar (NEMA 17)}
Cada bobina tiene dos terminales y la corriente debe invertirse para completar la
secuencia de pasos. Requiere un puente H por fase (integrado en \emph{drivers} como
A4988/DRV8825).
\begin{itemize}
  \item \textbf{Cables}: 4 (dos por bobina: $A+$, $A-$, $B+$, $B-$).
  \item \textbf{Voltaje}: \SIrange{12}{24}{\volt} típico.
  \item \textbf{Corriente por fase}: \SIrange{0.5}{2.5}{\ampere} según modelo.
  \item \textbf{Pasos por revolución}: 200 ($1.8\si{\degree}$/paso).
  \item \textbf{Torque}: \SIrange{2}{6}{\kilo\gram\centi\meter} (alto, adecuado para CNC,
        robótica).
  \item \textbf{Microstepping}: 1/2, 1/4, 1/8, 1/16 con A4988 (hasta 1/32 con DRV8825).
\end{itemize}

\subsection{Modos de excitación}
\subsubsection{Wave Drive (1 fase ON)}
Solo una bobina energizada a la vez. Consumo bajo, menor torque, mayor vibración.
Secuencia: A$\to$B$\to$C$\to$D$\to$A\ldots

\subsubsection{Full-Step (2 fases ON)}
Dos bobinas adyacentes energizadas simultáneamente. Máximo torque por paso. El rotor se
alinea entre dos polos. Para el 28BYJ-48 con reducción: \textbf{2048 pasos/revolución}.
\begin{table}[H]\centering
\caption{Secuencia \emph{full-step} para motor unipolar (4 fases).}
\begin{tabular}{ccccc}
\toprule
Paso & IN1 & IN2 & IN3 & IN4 \\
\midrule
1 & 1 & 0 & 0 & 0 \\
2 & 0 & 1 & 0 & 0 \\
3 & 0 & 0 & 1 & 0 \\
4 & 0 & 0 & 0 & 1 \\
\bottomrule
\end{tabular}
\end{table}

\subsubsection{Half-Step (8 fases)}
Alterna entre activar una y dos bobinas, duplicando la resolución angular. Para el
28BYJ-48 con reducción: \textbf{4096 pasos/revolución} ($\sim 0.088\si{\degree}$/paso).
Movimiento más suave, menor vibración, torque ligeramente inferior al \emph{full-step} de
2 fases.
\begin{table}[H]\centering
\caption{Secuencia \emph{half-step} para motor unipolar (8 fases).}
\begin{tabular}{ccccl}
\toprule
Paso & IN1 & IN2 & IN3 & IN4 \quad Descripción \\
\midrule
1 & 1 & 0 & 0 & 0 \quad Fase A \\
2 & 1 & 1 & 0 & 0 \quad Fases A+B \\
3 & 0 & 1 & 0 & 0 \quad Fase B \\
4 & 0 & 1 & 1 & 0 \quad Fases B+C \\
5 & 0 & 0 & 1 & 0 \quad Fase C \\
6 & 0 & 0 & 1 & 1 \quad Fases C+D \\
7 & 0 & 0 & 0 & 1 \quad Fase D \\
8 & 1 & 0 & 0 & 1 \quad Fases D+A \\
\bottomrule
\end{tabular}
\end{table}

\subsection{Resolución angular y reducción mecánica}
El 28BYJ-48 contiene un tren de engranajes interno con relación $\approx 64{:}1$. El
ángulo de paso del rotor es $5.625\si{\degree}$, lo que da 64 pasos por revolución
interna.
\begin{align}
  \text{SPR}_\text{full-step} &= \frac{360\si{\degree}}{5.625\si{\degree}}\times 64 =
        64\times 32 = 2048 \text{ pasos/revolución del eje}, \\
  \text{SPR}_\text{half-step} &= 2\times\text{SPR}_\text{full-step} = 4096
        \text{ pasos/revolución del eje}, \\
  \theta_\text{paso,half} &= \frac{360\si{\degree}}{4096}\approx 0.088\si{\degree}/\text{paso}.
\end{align}
\begin{table}[H]\centering
\caption{Resolución angular según modo de excitación (28BYJ-48).}
\begin{tabular}{lccc}
\toprule
Modo de excitación & Pasos/rev (eje) & Ángulo/paso & Resolución \\
\midrule
Full-step (4 fases) & 2048 & $0.176\si{\degree}$ & Media \\
Half-step (8 fases) & 4096 & $0.088\si{\degree}$ & Alta \\
\bottomrule
\end{tabular}
\end{table}

\subsection{Relación RPM — intervalo entre pasos}
La velocidad angular se controla mediante el retardo entre pasos consecutivos:
\begin{align}
  f_\text{paso} &= \frac{\text{RPM}\times\text{SPR}}{60} \quad [\text{pasos/s}], \\
  T_\text{paso} &= \frac{1}{f_\text{paso}} = \frac{60}{\text{RPM}\times\text{SPR}} \quad [\text{s}].
\end{align}
Por ejemplo, para 10\,RPM en modo \emph{half-step}:
$T_\text{paso} = 60/(10\times 4096) = \SI{1.46}{\milli\second}$/paso. El 28BYJ-48 tiene un
retardo mínimo de $\sim\SI{2}{\milli\second}$ entre pasos antes de perder la
sincronización, lo que limita la velocidad máxima a $\sim$15\,RPM en \emph{half-step}.

\subsection{El driver ULN2003}
El ULN2003A es un arreglo de \textbf{7 pares Darlington NPN} con diodos de protección
internos (\emph{flyback}) integrados. Es el complemento natural del motor unipolar, ya
que cada par Darlington actúa como interruptor de baja impedancia para una bobina.
\begin{table}[H]\centering
\caption{Características eléctricas del ULN2003A.}
\begin{tabular}{lll}
\toprule
Parámetro & Valor & Observación \\
\midrule
Canales & 7 (pares Darlington) & Se usan 4 para el 28BYJ-48 \\
$V_{CE(\text{sat})}$ & $\le\SI{1.6}{\volt}$ & A $I_C=\SI{350}{\milli\ampere}$ \\
$I_{C,\max}$ & 500\,mA por canal & Suficiente para el 28BYJ-48 ($\sim$200\,mA) \\
$V_{I(\text{on})}$ & $\ge\SI{2.4}{\volt}$ & Compatible con 3.3\,V y 5\,V \\
Diodos \emph{flyback} & Integrados (\emph{clamp}) & Protección contra transitorios inductivos \\
$V_\text{clamp}$ & $\le\SI{50}{\volt}$ & Adecuado para motores de 5--12\,V \\
\bottomrule
\end{tabular}
\end{table}

\begin{nota}[Diodos de protección integrados]
A diferencia de la P6, donde se seleccionó explícitamente un diodo \emph{flyback}
externo para cada carga inductiva, el ULN2003 \textbf{incluye diodos de \emph{clamp}
internos} entre cada salida y el pin COM (pin 9). Este pin debe conectarse a $V_{CC}$
del motor para que los diodos proporcionen el camino de recirculación de corriente al
apagar cada bobina. Si el pin COM se deja flotante, la protección se pierde y los
transitorios inductivos pueden dañar los Darlington.
\end{nota}

\subsection{Control STEP/DIR con A4988 (alternativa)}
El A4988 es un \emph{driver} de \emph{microstepping} con traductor incorporado para
motores bipolares. Su interfaz simplificada requiere solo dos señales:
\textbf{STEP} (un pulso = un micropaso) y \textbf{DIR} (dirección).
\begin{table}[H]\centering
\caption{Comparación de drivers: ULN2003 vs A4988.}
\begin{tabular}{lll}
\toprule
Característica & ULN2003 + 28BYJ-48 & A4988 + NEMA 17 \\
\midrule
Tipo de motor & Unipolar (5 cables) & Bipolar (4 cables) \\
Control MCU & 4 pines (secuencia) & 2 pines (STEP + DIR) \\
Secuencia en SW & Sí (tabla en firmware) & No (traductor interno) \\
Microstepping & No nativo & 1/2, 1/4, 1/8, 1/16 \\
Voltaje motor & 5\,V & 12--24\,V \\
Corriente por fase & $\sim$200\,mA & 0.5--2.5\,A \\
Torque & Bajo ($\sim$0.4\,kg$\cdot$cm) & Alto ($\sim$4\,kg$\cdot$cm) \\
Velocidad máx & $\sim$15\,RPM & $\sim$600\,RPM \\
Protección \emph{flyback} & Integrada (\emph{clamp}) & Puente H interno \\
Fuente externa & Opcional (VSYS ok) & Obligatoria \\
Costo & Bajo & Medio \\
Aplicación & Indicadores, prototipo & CNC, robótica, actuación \\
\bottomrule
\end{tabular}
\end{table}

\subsection{Aplicación en instrumentación aeronáutica: ARINC 429 Label 204}
En los instrumentos analógicos-electrónicos de cabina (\emph{glass-analog hybrid}), la
aguja es movida por un motor a pasos. El sistema recibe una palabra ARINC 429 con un
Label específico (por ejemplo, Label 204 para corrección barométrica o Label 203 para
indicador de RPM). El procesador decodifica el valor BNR (\emph{Binary Number
Representation}), calcula la posición angular correspondiente y envía los pasos
necesarios al \emph{driver} para actualizar la carátula.
\begin{table}[H]\centering
\caption{Estructura de una palabra ARINC 429 BNR (Label 204, corrección barométrica).}
\begin{tabular}{llll}
\toprule
Bits & Campo & Contenido & Ejemplo \\
\midrule
1--8 & Label & 204 (octal, invertido) & 00110100 \\
9--10 & SDI & Source ID & 00 \\
11--28 & Datos BNR & Ángulo codificado ($0$--$360\si{\degree}$) & Variable \\
29 & Signo/SSM & Estado (00 = Normal) & 0 \\
30--31 & SSM & Sign/Status Matrix & 00 \\
32 & P & Paridad impar & Calculada \\
\bottomrule
\end{tabular}
\end{table}
El sistema de posicionamiento debe: (1) validar la paridad impar (bit 32) para detectar
errores de transmisión; (2) verificar el SSM (bits 30--31) para confirmar datos válidos
(00 = \emph{Normal Operation}); (3) extraer el ángulo BNR del campo de datos; (4)
calcular la diferencia entre la posición actual y la deseada ($\Delta$ pasos); (5) elegir
la dirección de menor recorrido para minimizar el tiempo de actualización.

% ============================================================================
\section{Consideraciones de Seguridad}
\begin{enumerate}
  \item \textbf{Límites de corriente}: los motores a pasos consumen corriente constante
        incluso cuando están detenidos (\emph{holding torque}). Si se mantienen
        energizados por períodos prolongados, las bobinas y el \emph{driver} se
        calientan. Llamar \code{release()} para desactivar las bobinas cuando no se
        necesita mantener posición.
  \item \textbf{Desacoplo de potencia}: instalar un condensador electrolítico de
        $100\,\mu$F / 16\,V cerca de los pines de alimentación del ULN2003 (o en
        VMOT--GND del A4988) para absorber transitorios de conmutación.
  \item \textbf{No desconectar bobinas con \emph{driver} energizado}: especialmente con
        el A4988, desconectar las bobinas del motor mientras el \emph{driver} suministra
        corriente genera un arco que destruye los transistores del puente H.
  \item \textbf{GND común obligatorio}: la tierra del MCU y la tierra de la fuente del
        motor deben estar conectadas. Sin GND común, las señales lógicas no tienen
        referencia y el \emph{driver} no responde correctamente.
  \item \textbf{Fuente adecuada}: alimentar el 28BYJ-48 con 5\,V, nunca con 12\,V
        (destruye las bobinas). Para NEMA 17, usar fuente externa de 12--24\,V con
        capacidad de corriente nominal del motor.
  \item \textbf{Fin de carrera para \emph{homing}}: si se implementa \emph{homing},
        verificar que el microswitch esté correctamente cableado y probado antes de
        ejecutar. Establecer siempre un límite máximo de pasos como seguridad contra
        falla del sensor.
\end{enumerate}

% ============================================================================
\section{Equipo y Materiales}
\begin{table}[H]\centering
\caption{Lista de componentes para la Práctica 7.}
\begin{tabular}{clll}
\toprule
Cant.\ & Componente & Especificación & Uso \\
\midrule
1 & Microcontrolador & ESP32 DevKit / RP2040 & Control de secuencias GPIO \\
1 & Cable USB & Micro-B / USB-C & Programación y alimentación \\
1 & Protoboard & 830 puntos & Prototipado \\
--- & Alambre AWG 22 & Varios colores & Conexiones \\
1 & 28BYJ-48 & Motor unipolar, 5\,V & Caso principal \\
1 & ULN2003 (módulo) & Arreglo Darlington & \emph{Driver} del 28BYJ-48 \\
1 & NEMA 17 (opcional) & Motor bipolar, $1.8\si{\degree}$/paso & Caso alternativo \\
1 & A4988 (opcional) & \emph{Driver} \emph{microstepping} & \emph{Driver} del NEMA 17 \\
1 & Microswitch & Normalmente abierto (NO) & Fin de carrera (\emph{homing}) \\
1 & Condensador & $100\,\mu$F / 16\,V & Desacoplo fuente motor \\
\bottomrule
\end{tabular}
\end{table}
\textbf{Equipo de laboratorio}: laptop con IDE (Thonny / VS Code + Pymakr) y monitor
serie; fuente de alimentación DC de 5\,V (para 28BYJ-48) o 12\,V (para NEMA 17);
multímetro digital (DMM) para medir corriente de bobina y voltajes; osciloscopio
(recomendado) para visualizar secuencias IN1--IN4 o pulsos STEP/DIR.

% ============================================================================
\section{Etapa de Diseño}

\subsection{Caso de estudio 1: secuencia de excitación (28BYJ-48 + ULN2003)}
\textbf{Objetivo}: verificar la correcta implementación de las secuencias
\emph{full-step} y \emph{half-step}, observar el movimiento del motor y medir con
osciloscopio las señales en IN1--IN4.

\textbf{Cálculos de resolución} (modo \emph{half-step}, el modo por defecto del
firmware): $\text{SPR}=4096$ pasos/revolución;
$\theta_\text{paso}=360/4096=0.0879\si{\degree}$/paso; pasos para $90\si{\degree}=1024$;
pasos para $180\si{\degree}=2048$; pasos para $360\si{\degree}=4096$.

\textbf{Cálculo de velocidad} (RPM$=10$, conservadora):
$T_\text{paso}=60/(10\times 4096)=\SI{1.46}{\milli\second}$/paso. Retardo mínimo
recomendado para el 28BYJ-48: $\sim\SI{2}{\milli\second}$. A \SI{2}{\milli\second} por
paso: RPM$=60/(0.002\times 4096)\approx 7.3$.

\textbf{Corriente y potencia del \emph{driver}}: corriente por fase del 28BYJ-48
$\sim\SI{200}{\milli\ampere}$. En modo \emph{half-step}, 1 o 2 fases activas
simultáneamente: $I_{1\,\text{fase}}=\SI{200}{\milli\ampere}$,
$I_{2\,\text{fases}}=2\times 200=\SI{400}{\milli\ampere}$,
$P_\text{motor,max}=V_{CC}\times I_{2\,\text{fases}}=5\times 0.4=\SI{2.0}{\watt}$. Caída
en el ULN2003 por canal: $V_{CE(\text{sat})}\le\SI{1.6}{\volt}$ a \SI{350}{\milli\ampere}.
Voltaje real en la bobina: $V_\text{bobina}=5-1.6=\SI{3.4}{\volt}$.

\begin{nota}[Caída $V_{CE}$ del Darlington]
La caída de $\sim$1.6\,V del par Darlington es significativa comparada con la
alimentación de 5\,V. Reduce el voltaje efectivo en la bobina al 68\,\%, lo que
disminuye ligeramente el torque comparado con la alimentación directa. Para el
28BYJ-48, esto es aceptable porque el motor fue diseñado para operar con esta caída.
\end{nota}

\subsection{Caso de estudio 2: homing — posición de referencia}
\textbf{Objetivo}: implementar un algoritmo de referenciación que gire el motor en
sentido antihorario hasta activar un microswitch (fin de carrera), estableciendo la
posición cero del instrumento.

\textbf{Diseño del algoritmo}: (1) reducir velocidad al 50\,\% de la nominal (mayor
fiabilidad cerca del sensor); (2) girar en sentido CCW (antihorario) paso a paso; (3) en
cada paso, leer el pin del \emph{endstop} (activo LOW con \emph{pull-up} interno); (4) al
detectar \emph{endstop}$=0$: detener, avanzar 10 pasos (CW) para liberar el sensor; (5)
establecer \code{posicion\_actual = 0}; (6) límite de seguridad: máximo 1000 pasos sin
detección $\Rightarrow$ abortar con error.

\subsection{Asignación de pines (A4988, alternativa)}
\begin{table}[H]\centering
\caption{Asignación de pines para A4988 según plataforma (alternativa).}
\begin{tabular}{llll}
\toprule
Señal & ESP32 & RP2040 & Descripción \\
\midrule
STEP & GPIO18 & GP18 & Pulso por cada micropaso \\
DIR & GPIO19 & GP19 & Dirección (HIGH=CW, LOW=CCW) \\
ENABLE & GPIO5 & GP5 & Habilitación (activo LOW) \\
ENDSTOP & GPIO4 & GP4 & Fin de carrera (\emph{pull-up}, a GND) \\
GND & GND & GND & Tierra común \\
VMOT & 12\,V ext.\ & 12\,V ext.\ & Alimentación NEMA 17 \\
\bottomrule
\end{tabular}
\end{table}

% ============================================================================
\section{Procedimiento Experimental}

\subsection{Paso 1: preparación del entorno}
\begin{enumerate}
  \item Instalar MicroPython en la placa (o configurar PlatformIO con
        \code{PRACTICE=7}).
  \item Verificar que todos los componentes estén disponibles.
  \item Con la fuente apagada, conectar el módulo ULN2003 al motor 28BYJ-48 (conector
        de 5 pines; verificar orientación).
  \item Conectar IN1--IN4 del ULN2003 a los GPIO correspondientes.
  \item Conectar VCC del módulo a 5\,V y GND a tierra común con el MCU.
  \item Si se usa \emph{homing}: conectar microswitch entre ENDSTOP y GND.
  \item Sincronizar archivos del repositorio: \code{main.py}, \code{boot.py},
        \code{lib/stepper\_uln2003.py}.
\end{enumerate}

\subsection{Modos de operación del firmware}
\begin{table}[H]\centering
\caption{Modos de operación del firmware P7.}
\begin{tabular}{cll}
\toprule
Modo & Nombre & Descripción \\
\midrule
1 & Jog manual & \code{+} avanza 1 paso, \code{-} retrocede 1 paso \\
2 & Mover N pasos & Ingresa pasos (positivo/negativo) y RPM; movimiento preciso \\
3 & Barrido & Avanza hasta fin de carrera o límite, retrocede, repite \\
4 & Homing & Retrocede hasta fin de carrera, libera, establece posición cero \\
5 & Info del driver & Muestra configuración: tipo, pines, SPR, RPM \\
6 & Registro CSV & Jitter de intervalo STEP; CSV \code{t\_us,dt\_us} \\
\bottomrule
\end{tabular}
\end{table}
Tecla \code{m} $\to$ regresar al menú en cualquier modo.

\paragraph{Modo 6 — recolección de datos y vista en vivo.} Ejecuta 300 pasos a
RPM constante midiendo el intervalo real entre cada uno con \code{ticks\_us}
(MicroPython) / \code{micros()} (C++). Alimenta la toolkit compartida
\code{tools/sisela\_signal/} (disponible en MicroPython y C++, mismo formato CSV):
\begin{lstlisting}[language=bash]
python -m sisela_signal capture --port COM5 --menu 6 --out cap.csv
python -m sisela_signal characterize --file cap.csv --col dt_us --mode adc
python -m sisela_signal live --port COM5 --menu 6 --cols dt_us
\end{lstlisting}

\subsection{Caso 1: verificación de secuencia de fases}
\textbf{Verificar con osciloscopio}: canales en IN1--IN4 con la secuencia
\emph{half-step} de 8 estados; intervalo entre cambios de estado $\sim$\SI{3}{\milli\second}
por defecto; los LEDs del módulo ULN2003 deben encenderse en secuencia. Si el motor
vibra pero no gira: intercambiar IN2 e IN3 en el cableado.

\subsection{Caso 2: homing (referenciación)}
\textbf{Verificar}: el motor retrocede hasta activar el microswitch; se detiene y avanza
10 pasos para liberar el sensor; si el \emph{endstop} no se detecta en 1000 pasos, aborta
con mensaje de error.

\subsection{Caso 3 (ejercicio de diseño): posicionamiento por ARINC 429 Label 204}
\begin{nota}
El firmware actual del repositorio (\code{main.py}, \code{p7.cpp}) \textbf{no}
implementa la recepción ni decodificación de tramas ARINC 429 — los modos 1--5
(Tabla anterior) cubren jog, N pasos, barrido, homing e info del driver, todos por
comandos de texto simple sobre el puerto serie. Lo que sigue es un
\textbf{ejercicio de diseño/extensión} para practicar la integración con el
protocolo ya estudiado en las Prácticas 2 y 8, no una funcionalidad presente en el
código entregado.
\end{nota}
Como extensión, el firmware podría recibir una trama hexadecimal (simulando Label 204),
validar la paridad impar, extraer el ángulo BNR (bits 11--28), calcular la trayectoria
mínima ($\Delta$ pasos con signo, plegando sobre $\pm\text{SPR}/2$) y mover el motor en
la dirección correspondiente.

\subsubsection{Ejemplo de cálculo}
Posición actual $90\si{\degree}$ (1024 pasos). Ángulo objetivo (de ARINC 429):
$270\si{\degree}$ (3072 pasos).
$$\Delta = 3072-1024 = 2048 \text{ pasos (CW)}.$$
$2048=\text{SPR}/2$, por lo que cualquier dirección es equivalente: 2048 pasos. Si el
objetivo fuera $300\si{\degree}$ (3413 pasos):
$$\Delta = 3413-1024=2389 > 2048 \Rightarrow \Delta_\text{adj}=2389-4096=-1707
\text{ (CCW, más corto)}.$$

\subsection{Caso 4: barrido automático (Modo 3)}
Avanza hasta el fin de carrera (o un límite de 400 pasos sin \emph{endstop}), retrocede,
y repite. Verificar que el retorno llega a la posición inicial dentro del error de
repetibilidad esperado.

\subsection{Implementación C++ (PlatformIO)}
El archivo \code{p7.cpp} implementa los modos usando la clase \code{LandingGear}
(\code{common/landing\_gear.h}), que encapsula tanto la secuencia ULN2003 como el control
STEP/DIR del A4988 según la macro \code{STEPPER\_ULN2003}, incluido el Modo 5 de
registro CSV (\code{common/siglab.h}) — mismo formato \code{t\_us,dt\_us} que la
versión MicroPython. Compilar con \code{pio run -e esp32dev} o \code{pio run -e
pico} (con \code{-DPRACTICE=7}).

% ============================================================================
\section{Alternativa de Trabajo en Casa (Multímetro + Simulación)}
\label{sec:casa}

Esta práctica es, junto con la P4, de las más viables para trabajar sin laboratorio: el
kit 28BYJ-48 + ULN2003 es económico ($\sim$US\$3--5) y el osciloscopio es
\textbf{recomendado, no obligatorio} — la verificación principal (¿gira el motor en el
sentido y ángulo correctos?) puede hacerse a simple vista con un transportador.

\paragraph{Con multímetro (en casa).} Medir la corriente de una bobina insertando el
amperímetro en serie entre una salida del ULN2003 y la bobina (motor detenido, modo
\emph{holding}): esperado $\sim$\SIrange{150}{200}{\milli\ampere} por fase activa.
Verificar continuidad de las 4 (ULN2003) o 2 (A4988) bobinas con el óhmetro. Verificar
5\,V (o 12\,V para NEMA 17) en la alimentación del \emph{driver}.

\paragraph{Con simulación.} Proteus incluye modelos razonablemente completos de motor a
pasos y de los \emph{drivers} ULN2003/A4988 — se puede simular la secuencia de excitación
y observarla en un osciloscopio virtual, e incluso animar el giro del motor virtual antes
de tener el motor físico. Es también una buena forma de practicar el cálculo de
trayectoria mínima del Caso 3 sin hardware.

\paragraph{Requiere laboratorio presencial.} La medición \emph{precisa} de repetibilidad
(Tabla~13, 10 repeticiones con error en grados) requiere el motor físico y, para
verificar el intervalo real entre pasos en microsegundos, un osciloscopio; a simple vista
con un transportador solo se puede estimar el error grueso ($\pm$varios grados), no la
precisión submilimétrica que exige el análisis de la \S11.

\begin{table}[H]\centering
\caption{Alternativas de trabajo en casa.}
\begin{tabular}{lccc}
\toprule
Actividad & Multímetro & Simulación & Requiere laboratorio \\
\midrule
Secuencia de fases (Caso 1) & --- & Sí (visual) & Recomendado (osciloscopio) \\
Corriente de bobina & Sí & --- & No \\
Homing (Caso 2) & --- & Sí & No (con el motor físico en casa) \\
Trayectoria mínima (Caso 3) & --- & Sí (lógica) & No \\
Repetibilidad de precisión (\S11) & --- & No & Sí, para el error en grados \\
\bottomrule
\end{tabular}
\end{table}

% ============================================================================
\section{Herramientas de Verificación}
\subsection{Verificación con osciloscopio}
\begin{enumerate}
  \item \textbf{Secuencia IN1--IN4 (ULN2003)}: conectar 4 canales a IN1--IN4. Observar la
        secuencia \emph{half-step} (8 estados) o \emph{full-step} (4 estados). Verificar
        el intervalo entre cambios de estado ($\sim$\SIrange{2}{5}{\milli\second}).
  \item \textbf{Señal STEP (A4988)}: tren de pulsos rectangulares. Frecuencia
        $=\text{RPM}\times\text{SPR}/60$. Para 60\,RPM con 200 pasos/rev:
        $f=\SI{200}{\hertz}$ ($T=\SI{5}{\milli\second}$).
  \item \textbf{Señal DIR (A4988)}: nivel constante durante movimiento en una dirección.
        Cambio de nivel = cambio de dirección.
  \item \textbf{Corriente de bobina}: colocar resistencia \emph{shunt} ($1\,\Omega$) en
        serie con una salida del ULN2003. Medir voltaje como $I=V_\text{shunt}/R_\text{shunt}$.
        Esperado: $\sim$\SIrange{150}{200}{\milli\ampere} por bobina para el 28BYJ-48.
\end{enumerate}
\subsection{Verificación con multímetro}
Medir voltaje en VCC del módulo ULN2003 (debe ser $\sim$5\,V); medir corriente total del
motor detenido (\emph{holding}): $\sim$\SIrange{200}{400}{\milli\ampere} (1--2 fases
activas); verificar continuidad de GND común entre MCU y módulo ULN2003.

% ============================================================================
\section{Registro de Datos y Análisis}

\subsection{Tabla de precisión de posicionamiento}
\begin{table}[H]\centering
\caption{Bitácora de precisión y desempeño del motor a pasos.}
\begin{tabular}{ccccc}
\toprule
Ángulo objetivo ($\si{\degree}$) & Pasos teóricos & Pasos ejecutados & Error medido ($\si{\degree}$) & Tiempo (s) \\
\midrule
0 (Home) & 0 & & & \\
90 & 1024 & & & \\
180 & 2048 & & & \\
270 & 3072 & & & \\
360 & 4096 & & & \\
\bottomrule
\end{tabular}
\end{table}

\subsection{Tabla de repetibilidad}
\begin{table}[H]\centering
\caption{Análisis de repetibilidad (10 repeticiones de $0\si{\degree}\to 180\si{\degree}$).}
\begin{tabular}{cccc}
\toprule
Repetición & Ángulo final ($\si{\degree}$) & Error ($\si{\degree}$) & Observaciones \\
\midrule
1--10 & & & \\
\midrule
Promedio & & & \\
Desv.\ estándar & & & \\
\bottomrule
\end{tabular}
\end{table}

\subsection{Tabla de velocidad y consumo}
\begin{table}[H]\centering
\caption{Consumo de corriente según modo de operación.}
\begin{tabular}{lccl}
\toprule
Estado & Corriente (mA) & Voltaje bobina (V) & Observaciones \\
\midrule
Idle (\code{release()}) & $\sim 0$ & & Bobinas desactivadas \\
Holding (1 fase) & $\sim 200$ & & Posición mantenida \\
Holding (2 fases) & $\sim 400$ & & Máximo torque \\
Movimiento (\emph{half-step}) & Variable & & Promedio medido \\
\bottomrule
\end{tabular}
\end{table}

\subsection{Gráficas requeridas}
\begin{enumerate}
  \item \textbf{Secuencia IN1--IN4}: oscilograma de las 4 señales mostrando la secuencia
        \emph{half-step} completa (8 estados).
  \item \textbf{Error angular vs ángulo comandado}: gráfica de dispersión con barras de
        error para la Tabla de precisión.
  \item \textbf{Histograma de repetibilidad}: distribución de errores de las 10
        repeticiones.
  \item \textbf{Corriente vs tiempo}: forma de onda de corriente durante movimiento
        (resistencia \emph{shunt} + osciloscopio).
\end{enumerate}

% ============================================================================
\section{Análisis de Resultados}
\begin{enumerate}
  \item \textbf{Análisis de repetibilidad}: repetir 10 veces el movimiento de
        $0\si{\degree}$ a $180\si{\degree}$. Medir el error de llegada final en cada
        intento. Calcular media y desviación estándar. Discutir si existe pérdida de
        pasos por resonancia o exceso de velocidad.
  \item \textbf{Full-step vs half-step}: comparar cualitativamente el ruido, vibración y
        suavidad del movimiento en ambos modos. Medir la resolución angular real y
        compararla con la teórica.
  \item \textbf{Efecto de la velocidad}: encontrar la velocidad máxima (RPM) a la que el
        motor opera sin perder pasos. Graficar error vs RPM. Discutir la relación
        torque-velocidad.
  \item \textbf{Validación de trayectoria mínima}: para el Caso 3, verificar que el motor
        elija la dirección de menor recorrido al cambiar entre posiciones distantes (ej.
        $90\si{\degree}\to 300\si{\degree}$).
  \item \textbf{Homing}: verificar que el procedimiento de referencia produce resultados
        consistentes. ¿Cuál es la variabilidad de la posición cero después de 5
        procedimientos de \emph{homing} consecutivos?
  \item \textbf{Consumo energético}: comparar el consumo en \emph{holding} vs
        movimiento. ¿Cuánta energía se ahorra al llamar \code{release()} cuando el motor
        está detenido? Discutir el compromiso entre consumo y estabilidad de posición
        (como un altavoz aeronave).
\end{enumerate}

% ============================================================================
\section{Preguntas de Discusión y Refuerzo}
\begin{enumerate}
  \item \textbf{Resonancia}: explique el fenómeno de resonancia en motores a pasos y
        cómo una rampa de aceleración por software ayuda a mitigar este problema en
        instrumentos de vuelo. ¿A qué frecuencia de paso ocurre típicamente la
        resonancia en el 28BYJ-48?
  \item \textbf{Torque de mantenimiento}: ¿qué sucede con el consumo de potencia si se
        dejan las bobinas energizadas mientras la aguja no se mueve? ¿Por qué se
        prefiere este consumo en lugar de desactivar el motor en un ambiente de
        vibración aeronáutica?
  \item \textbf{Resolución}: si el instrumento requiere una resolución de
        $0.01\si{\degree}$, ¿cómo modificaría el hardware (engranajes adicionales) o el
        software (\emph{microstepping} con A4988 a 1/16) para lograrlo? Calcule los
        pasos por revolución necesarios.
  \item \textbf{Detección de fallas}: ¿cómo puede el software detectar que el motor ha
        perdido pasos sin un \emph{encoder}? Proponga un método de verificación
        periódica contra el sensor de \emph{homing}.
  \item \textbf{ULN2003 vs puente H}: ¿por qué el ULN2003 no puede controlar un motor
        bipolar (NEMA 17)? Explique la necesidad de invertir la corriente y cómo lo
        resuelve el A4988.
  \item \textbf{DO-178C}: en el contexto de instrumentos certificados, ¿a qué nivel DAL
        (\emph{Design Assurance Level}) clasificaría el software de un indicador de
        altitud? ¿Qué implicaciones tiene para las pruebas y verificación del código?
  \item \textbf{Pérdida de pasos bajo carga}: si se aplica una carga mecánica progresiva
        al eje del motor, ¿cómo se manifiesta la pérdida de pasos? ¿El motor se detiene
        abruptamente o acumula error gradualmente?
  \item \textbf{Consumo en \emph{standby}}: en un panel de instrumentos con 6
        indicadores de motor a pasos, estime el consumo total si cada motor mantiene las
        bobinas energizadas continuamente. ¿Es aceptable para un sistema alimentado por
        batería de emergencia?
  \item \textbf{Comparación servo vs \emph{stepper}}: compare las ventajas y
        desventajas de un servomotor (P5) versus un motor a pasos (P7) para posicionar
        una aguja de instrumento. Considere resolución, costo, complejidad del
        \emph{driver} y fiabilidad.
\end{enumerate}

% ============================================================================
\section{Entregables (Formato AIAA)}
\begin{enumerate}
  \item \textbf{Código fuente}: \code{main.py} + \code{lib/stepper\_uln2003.py}
        (MicroPython) o \code{p7.cpp} (C++).
  \item \textbf{Memoria de cálculo}: SPR, resolución angular, intervalos de paso,
        corriente de bobina.
  \item \textbf{Tabla de secuencias}: \emph{full-step} y \emph{half-step} implementadas,
        con verificación por osciloscopio.
  \item \textbf{Tabla de precisión} (\S10.1): ángulos objetivo vs medidos.
  \item \textbf{Análisis de repetibilidad} (\S10.2): media, desviación estándar,
        histograma.
  \item \textbf{Oscilogramas}: secuencia IN1--IN4 (o STEP/DIR) capturada
        experimentalmente.
  \item \textbf{Análisis}: mínimo 4 de los 6 puntos de la \S11.
  \item \textbf{Discusión}: responder $\ge 4$ de las 9 preguntas.
  \item \textbf{Conclusiones}: hallazgos, fuentes de error, mejoras propuestas.
  \item \textbf{Referencias}: $\ge 3$ fuentes (\emph{datasheets} del 28BYJ-48 y ULN2003A,
        libro de actuadores, normativa DO-160G).
\end{enumerate}

% ============================================================================
\section{Referencias}
\begin{enumerate}
  \item Kenjo, T., Sugawara, A. \emph{Stepping Motors and Their Microprocessor
        Controls}, 2nd ed., Oxford University Press, 1994.
  \item Texas Instruments. \emph{ULN2003A Seven Darlington Arrays Datasheet}, SLRS027,
        Rev.\ T.
  \item Kiatronics. \emph{28BYJ-48 --- 5V Stepper Motor Technical Specifications}, 2015.
  \item Allegro MicroSystems. \emph{A4988 DMOS Microstepping Driver with Translator
        Datasheet}, Rev.\ 5.
  \item ARINC. \emph{Specification 429 Part 1}: Mark 33 Digital Information Transfer
        System (DITS), 2004.
  \item RTCA. \emph{DO-160G}: Environmental Conditions and Test Procedures for Airborne
        Equipment, Sec.\ 16 (Power Input), 2010.
  \item RTCA. \emph{DO-178C}: Software Considerations in Airborne Systems and Equipment
        Certification, 2011.
  \item Moir, I., Seabridge, A. \emph{Aircraft Systems: Mechanical, Electrical, and
        Avionics Subsystems Integration}, 3rd ed., Wiley, 2008.
  \item Collins Aerospace. \emph{Pro Line Fusion --- Air Data and Attitude Heading
        Reference System}, Technical Description, 2020.
  \item MicroPython. \emph{machine.Pin Documentation}.
\end{enumerate}

% ============================================================================
\appendix
\section{Código Fuente del Repositorio}
\begin{table}[H]\centering
\caption{Estructura de archivos de la Práctica 7.}
\small
\begin{tabular}{p{3.4cm}p{9.8cm}}
\toprule
Plataforma & Archivo(s) / Descripción \\
\midrule
MicroPython/ESP32 & \code{P7/main.py} (5 modos), \code{P7/boot.py},
                    \code{P7/lib/stepper\_uln2003.py}, \code{P7/lib/stepper\_a4988.py} \\
MicroPython/RP2040 & Adaptación de pines (idéntica lógica), mismos \emph{drivers} \\
C++ / PlatformIO & \code{practices/p7.cpp}, \code{common/landing\_gear.h}
                   (\code{LandingGear}, control ULN2003/A4988) \\
\bottomrule
\end{tabular}
\end{table}

\subsection{Diferencias entre plataformas}
\begin{table}[H]\centering
\caption{Resumen de diferencias P7 entre plataformas.}
\begin{tabular}{lll}
\toprule
Aspecto & ESP32 & RP2040 \\
\midrule
Pines ULN2003 & GPIO26, 25, 33, 32 & GP26, 27, 28, 22 \\
Pines A4988 & GPIO18, 19, 5 & GP18, 19, 5 \\
Endstop & GPIO4 & GP4 \\
GPIO máx & 39 (no todos de salida) & 30 (todos bidireccionales) \\
\emph{Pull-up} interno & Sí (GPIOs de entrada) & Sí (todos los GPIOs) \\
Lógica firmware & Idéntica & Idéntica \\
\bottomrule
\end{tabular}
\end{table}

\section{Contexto Aeronáutico Complementario}
\subsection{Instrumentación analógica de cabina}
En \emph{cockpits} \emph{glass-analog hybrid}, los indicadores de aguja siguen siendo
comunes como respaldo o como instrumentos primarios en aviones de aviación general. Cada
instrumento contiene un micro motor a pasos (frecuentemente unipolar con reducción) que
mueve la aguja en respuesta a datos del bus ARINC 429. El proceso de operación es: (1) al
encender (\emph{power-up}), el instrumento ejecuta una secuencia de \emph{homing}: la
aguja barre toda la carátula hasta un sensor óptico o magnético que define el cero
mecánico; (2) el procesador se suscribe a las palabras ARINC 429 pertinentes; (3) al
recibir cada palabra válida, calcula la posición angular proporcional al valor BNR y
envía los pasos necesarios al \emph{driver}, eligiendo la trayectoria mínima; (4) si la
trama indica estado inválido o se pierde la señal por más de 2\,s, la aguja se mueve a
una posición \emph{flagged} (cruz roja) y se ilumina una bandera de falla.

\subsection{Ventajas de los motores a pasos en instrumentación}
\textbf{Posicionamiento absoluto sin \emph{encoder}}: reducción de peso, costo y
complejidad. \textbf{Alta resolución}: con reducción mecánica, se alcanzan resoluciones
de $<0.1\si{\degree}$. \textbf{Holding torque}: la aguja mantiene su posición bajo
vibración (frecuencias típicas en aeronaves: 10--500\,Hz según DO-160G Sección 8).
\textbf{Compatibilidad digital}: control directo desde el bus de datos sin conversiones
D/A.

\subsection{Normativa aplicable}
\textbf{DO-160G, Sección 8} (Vibración): los motores a pasos deben mantener la precisión
de indicación bajo los perfiles de vibración especificados para la zona de montaje del
instrumento. \textbf{DO-160G, Sección 16} (Potencia de entrada): el instrumento debe
operar correctamente bajo variaciones de voltaje del bus DC (típicamente 18--32\,VDC
para bus de 28\,V).

\end{document}
